% ============================================
% PLANTILLA DE TESIS - LaTeX
% Incluye: índice, tablas, imágenes, anexos,
% enlaces externos, citas APA/IEEE
% ============================================

\documentclass[12pt, a4paper]{report}

% ---- CODIFICACIÓN Y IDIOMA ----
\usepackage[utf8]{inputenc}
\usepackage[T1]{fontenc}
\usepackage[spanish]{babel}

% ---- MÁRGENES ----
\usepackage[margin=2.5cm]{geometry}

% ---- IMÁGENES ----
\usepackage{graphicx}
\graphicspath{{imagenes/}}

% ---- TABLAS MEJORADAS ----
\usepackage{booktabs}
\usepackage{array}
\usepackage{caption}

% ---- ENLACES (internos y externos) ----
\usepackage{hyperref}
\hypersetup{
    colorlinks=true,
    linkcolor=blue,
    filecolor=magenta,
    urlcolor=cyan,
    citecolor=green,
    pdftitle={Mi Tesis},
    pdfauthor={Tu Nombre},
}

% ---- BIBLIOGRAFÍA (cambiar apa por ieee si prefieres) ----
\usepackage[backend=biber, style=apa, sorting=nyt]{biblatex}
% Para IEEE usar: style=ieee
\addbibresource{referencias.bib}

% ---- ANEXOS ----
\usepackage[toc,page]{appendix}

% ---- ÍNDICE DE FIGURAS Y TABLAS ----
\usepackage{tocloft}

% ---- CÓDIGO FUENTE (opcional) ----
\usepackage{listings}
\lstset{
    basicstyle=\ttfamily\small,
    breaklines=true,
    frame=single
}

% ============================================
% INICIO DEL DOCUMENTO
% ============================================
\begin{document}

% ---- PORTADA ----
\begin{titlepage}
    \centering
    \vspace*{2cm}
    {\Huge\bfseries Título de la Tesis\par}
    \vspace{1.5cm}
    {\Large Subtítulo o línea descriptiva\par}
    \vspace{2cm}
    {\large\itshape Tu Nombre Completo\par}
    \vspace{1cm}
    {\large Director: Dr. Nombre del Director\par}
    \vfill
    {\large Universidad / Institución\par}
    {\large Facultad / Departamento\par}
    \vspace{1cm}
    {\large Enero 2026\par}
\end{titlepage}

% ---- ÍNDICES ----
\tableofcontents          % Índice general
\listoffigures            % Índice de figuras
\listoftables             % Índice de tablas
\newpage

% ---- CAPÍTULOS ----
\chapter{Introducción}
\label{cap:introduccion}

Este es un ejemplo de tesis en \LaTeX. Aquí puedes escribir la introducción de tu trabajo.

\section{Contexto del problema}
El problema que aborda esta tesis es muy importante según \textcite{knuth1984}.

\section{Objetivos}
\subsection{Objetivo general}
Demostrar el uso de \LaTeX{} para tesis académicas.

\subsection{Objetivos específicos}
\begin{itemize}
    \item Aprender a usar tablas y figuras
    \item Dominar el sistema de citas
    \item Crear documentos profesionales
\end{itemize}

% ---- EJEMPLO DE ENLACE EXTERNO ----
Para más información, visita \href{https://www.latex-project.org/}{el sitio oficial de LaTeX}.

\chapter{Marco Teórico}
\label{cap:marco}

\section{Fundamentos}
Según \textcite{lamport1994}, \LaTeX{} es un sistema de preparación de documentos de alta calidad.

También se puede citar entre paréntesis \parencite{knuth1984}.

% ---- EJEMPLO DE FIGURA ----
\section{Figuras}
La Figura \ref{fig:ejemplo} muestra un diagrama de ejemplo.

\begin{figure}[htbp]
    \centering
    % Descomentar cuando tengas una imagen:
    % \includegraphics[width=0.7\textwidth]{mi_imagen.png}
    \fbox{\parbox{0.6\textwidth}{\centering\vspace{3cm}[Aquí va tu imagen]\\mi\_imagen.png\vspace{3cm}}}
    \caption{Descripción de la figura de ejemplo}
    \label{fig:ejemplo}
\end{figure}

% ---- EJEMPLO DE TABLA ----
\section{Tablas}
La Tabla \ref{tab:ejemplo} presenta datos de ejemplo.

\begin{table}[htbp]
    \centering
    \caption{Comparación de características}
    \label{tab:ejemplo}
    \begin{tabular}{lcc}
        \toprule
        \textbf{Característica} & \textbf{Opción A} & \textbf{Opción B} \\
        \midrule
        Velocidad & Alta & Media \\
        Costo & Bajo & Alto \\
        Facilidad & Media & Alta \\
        \bottomrule
    \end{tabular}
\end{table}

\chapter{Metodología}
\label{cap:metodologia}

Aquí describes tu metodología. Puedes referenciar capítulos anteriores como el Capítulo \ref{cap:introduccion}.

\chapter{Resultados}
\label{cap:resultados}

Presenta tus resultados aquí.

\chapter{Conclusiones}
\label{cap:conclusiones}

Tus conclusiones van aquí.

% ---- BIBLIOGRAFÍA ----
\printbibliography[heading=bibintoc, title={Referencias}]

% ---- ANEXOS ----
\begin{appendices}
\chapter{Datos adicionales}
\label{anexo:datos}

Aquí van tus anexos. Puedes incluir tablas extensas, código fuente, cuestionarios, etc.

\begin{lstlisting}[language=Python, caption={Código de ejemplo}]
def hola_mundo():
    print("Hola desde LaTeX!")
    return True
\end{lstlisting}

\chapter{Cuestionario utilizado}
\label{anexo:cuestionario}

El cuestionario completo se presenta a continuación...

\end{appendices}

\end{document}
